\section{The discrete coset-ring step}
\label{sec:s5-coset}

Write $\mathbb R_d$ for the additive real group with the discrete
topology. Its coset ring is the Boolean algebra generated by the
cosets of all its abstract subgroups. Local finiteness below always
refers to the usual topology of $\mathbb R$: every compact interval
meets the set in finitely many points. This distinction matters
because an abstract subgroup used in a coset-ring representation
can be dense in the usual topology.

\paragraph{MA-FQ-002 (local finiteness removes dense cosets).}
If $E$ belongs to the coset ring of $\mathbb R_d$ and is locally
finite in $\mathbb R$, there are a finite set $F$ and finitely many
arithmetic progressions $\beta_j+\alpha_j\mathbb Z$, with
$\alpha_j>0$, such that
\[
 E\mathbin{\triangle}\bigcup_{j=1}^{J}
                  (\beta_j+\alpha_j\mathbb Z)=F.
 \tag{MA-FQ.2}
\]
The empty union is allowed. Consequently, if an integer-valued
function $a\in B(\mathbb R_d)$ has locally finite support in
$\mathbb R$, then there are integers $c_j,d_\ell$ and real numbers
$\alpha_j>0$, $\beta_j$, and $\gamma_\ell$, with finite index sets,
such that
\[
 a(t)=\sum_{j=1}^{J}c_j
                 1_{\beta_j+\alpha_j\mathbb Z}(t)
       +\sum_{\ell=1}^{L}d_\ell 1_{\{\gamma_\ell\}}(t).
 \tag{MA-FQ.3}
\]
In particular its atomic measure is the corresponding finite
integer linear combination of translated lattice combs and point
masses. No uniform separation of distinct points is assumed.

\begin{proof}
We first prove the assertion about sets, including the issue of
Boolean differences of dense subgroups. Every element of the
coset ring is a finite union of atoms of the form
\[
 A\setminus\bigcup_{i=1}^{n}B_i,
 \tag{MA-FQ.4}
\]
where $A$ is a coset of a subgroup, and each $B_i$ is a coset
contained in $A$. Indeed, take the finite Boolean algebra
generated by the cosets in one expression for $E$, intersect the
positive cosets, and intersect each excluded coset with that
intersection. If there are no positive cosets, use
$A=\mathbb R$. Empty atoms may be discarded.

Suppose an atom in (MA-FQ.4) is nonempty and that the subgroup
underlying $A$ is dense. Translate a surviving point to $0$.
We then have $A=H$, with $H$ a dense subgroup of $\mathbb R$,
and $B_i=x_i+H_i\subset H$, with $0\notin B_i$.
Consider the finite family
\[
 \left\{H\cap\bigcap_{i\in I}H_i:
                         I\subseteq\{1,\ldots,n\}\right\}.
\]
Choose an inclusion-minimal member $K$ that is dense in
$\mathbb R$; such a member exists because $H$ is in the family.
For each $i$, if $K\cap H_i$ is dense, minimality gives
$K\cap H_i=K$. Thus $K\subset H_i$, and $K\cap B_i$ is
empty: an intersection would imply $B_i=H_i$, contrary to
$0\notin B_i$.

In the other case, $K\cap H_i$ has proper closed subgroup
closure in $\mathbb R$. Such a closure is either $\{0\}$ or a
lattice, and is nowhere dense. The intersection $K\cap B_i$
is either empty or a coset of $K\cap H_i$, so its closure is
also nowhere dense. A finite union of these closures cannot
fill any open interval. Every open interval therefore contains
a nonempty open subinterval avoiding them, and the density of
$K$ supplies a point of $K$ in that subinterval. It follows that
\[
 K\setminus\bigcup_{i=1}^{n}B_i
\]
is dense in $\mathbb R$. This set lies in the atom, so a
nonempty atom with dense ambient subgroup cannot be a subset
of a locally finite set. Every nonempty atom in the expression
for $E$ must therefore have discrete ambient subgroup.

A discrete subgroup of $\mathbb R$ is either $\{0\}$ or
$\alpha\mathbb Z$ for some $\alpha>0$. The first case gives a
singleton atom. In the second case identify $A$ with
$\mathbb Z$. Each excluded coset $B_i$ is then either a
singleton or a residue class modulo a positive integer $n_i$.
Let $N$ be the least common multiple of the finitely many
$n_i$, with $N=1$ when there are none. Removing the residue
classes leaves a union of residue classes modulo $N$; removing
the singleton exclusions changes it at only finitely many
points. Hence every atom is a finite union of arithmetic
progressions up to a finite symmetric difference. Taking the
finite union proves (MA-FQ.2).

For the function assertion, Sanders's Theorem~12.3 applies to
the locally compact Abelian group $\mathbb R_d$. It gives
finitely many subgroups $H_i$ and integer-valued functions
$z_i$ on $\mathbb R_d/H_i$ with finite $\ell^1$ norms, such
that
\[
 a=\sum_i\sum_{W\in\mathbb R_d/H_i}z_i(W)1_W.
 \tag{MA-FQ.5}
\]
An integer-valued $\ell^1$ function has finite support.
Thus (MA-FQ.5) uses only finitely many cosets. The function
$a$ is bounded, hence has finite integer range, and each level
set $E_r=\{t:a(t)=r\}$ lies in the finite Boolean algebra
generated by those cosets. For $r\ne0$, the set $E_r$ is a
subset of the locally finite support of $a$; apply (MA-FQ.2)
to it.

It remains to pass from a union of progressions to an integer
linear combination of their indicators. Inclusion--exclusion
does this: every nonempty intersection of finitely many
progressions is either another progression or a singleton,
because it is a coset of the intersection of their discrete
subgroups. The finite symmetric differences add only finitely
many singleton indicators. Sum the resulting expressions for
$r1_{E_r}$ over the finite nonzero range of $a$. This proves
(MA-FQ.3), including the assertion about the associated atomic
measure.
\end{proof}

\paragraph{Primary sources and exact scope.}
The set classification is Schreiber's Example~1.6, p.~808;
its dense-coset input is Lemma~1.4 and Corollary~1.5, p.~807.
The proof above supplies the needed real-line argument directly,
choosing an actual surviving point before taking intersections.
Kolountzakis gives a second general proof in
\S3.2, Lemma~2 and Theorem~3, author pp.~7--8 of
\path{arXiv:math/9904065v1}; the real-line classification is
display (9), p.~7, attributed there to Rosenthal. The
Fourier--Stieltjes input is Sanders's Theorem~12.3 and definition
(12.2); only its finite qualitative representation is used.
The present argument does not require its quantitative bounds
and does not reprove Cohen's idempotent theorem.

Schreiber, B. M. (1970). On the coset ring and strong Ditkin
sets. \emph{Pacific Journal of Mathematics, 32}(3), 805--812.
\url{https://doi.org/10.2140/pjm.1970.32.805}.

Kolountzakis, M. N. (2000). On the structure of multiple
translational tilings by polygonal regions.
\emph{Discrete \& Computational Geometry, 23}(4), 537--553.
\url{https://doi.org/10.1007/s004540010014}.

Sanders, T. (2020). Bounds in Cohen's idempotent theorem.
\emph{Journal of Fourier Analysis and Applications, 26},
Article 25. \url{https://doi.org/10.1007/s00041-020-09732-y}.
